\documentclass[]{article}

%opening
\title{Finding optimal solutions to a Teacher Assignment Problem using MILP- and SMT-Solvers}
\author{Benedikt Schenk}



\begin{document}

\maketitle

\begin{abstract}
	
	The area of Teacher Assignment Problems is a topic of high interest, not only for schools and universities. Therefore, over the years many different solution-strategies evolved, each trying to keep computation time low whilst still finding an optimum or near-optimum. We look at a specific real-world version of this problem: teachers have to be assigned to courses in geographically different locations according to their distance, preferences, as well as availability constraints. We reformulate the problem as a linear optimization problem over integer variables by linearizing the non-linear parts of the objective function. Two different approaches for the optimization are used and compared: Using the MILP-Solver Pulp and using the SMT-Solver Z3 with the vZ extension.

\end{abstract}

\pagebreak


\tableofcontents

\pagebreak

\section{Introduction}
\subsection{Overview over Teacher Assignment Problems}
\begin{itemize}
	\item what is a teacher assignment problem? 

	\item Transportation Problem: suppliers with supply, destinations with demand, edge between all suppliers and destinations, cost for each edge, decision variable for each edge
	\item Assignment Problem: any source (agent) can be assigned to any destination (task), each edge has a cost and a binary decision variable. Demand of each source is 0 or 1.
	\item Teacher Assignment Problem: TP, but demand can be bigger than 1, not every source is connected to every destination, additional objective regarding distribution of workdays and additional constraints
	\item origin of teacher assignment problems / typical use cases
\end{itemize}

\subsection{Related work}
\begin{itemize}
	\item overview of scientific material, different strategies, algorithms
\end{itemize}

\subsection{Problem description}
\begin{itemize}
	\item describing the concrete problem which is being solved in this thesis
	\item stress that also unsatisfiable scenarios should be considered
\end{itemize}

\section{Formulating the problem}
\begin{itemize}
	\item Linear Solvers as flexible, fast way to find perfect optimum $\rightarrow$ aiming to find linear cost function and linear constraints
\end{itemize}



\subsection{Cost function}
\begin{itemize}
	\item what needs to be considered?
	\item how are the different costs balanced?  normalization, possibility of custom weights
	\item problem of standard deviation (not linear) $\rightarrow$ optimization of average relative deviation and maximum relative deviation
	\item formula of the cost function
\end{itemize}

\subsection{Constraints}
\begin{itemize}
	\item formulas of the constraints
	\item explaining the workaround to get average relative deviations, as well as the maximum of them
\end{itemize}

\subsection{Balancing the weights}
\begin{itemize}
	\item How is standard deviation best approximated (in my use case)
	\item How should you weigh standard deviation against priority against distance
\end{itemize}

\section{Computing the optimum}
\begin{itemize}
	\item ILP Problem, therefore using the MILP Solver pulp makes sense
	\item vZ is tried, to see how it compares
\end{itemize}

\subsection{Comparison of the solvers}
\begin{itemize}
	\item What goes on inside of PuLP?
	\item PuLP easy to handle, no issues
	\item What goes on inside of vZ?
	\item vZ has issues of mixing Ints and Reals in the objective function
	\item a custom binary search was implemented for Z3 in order to create custom cutoff-point
\end{itemize}

\subsection{Considering unsatisfiability}
\begin{itemize}
	\item use Z3 to find out which teachers have to be available additionaly for which events, in order to make it satisfiable
	\item the problem is solved with less strict constraints regarding the needed number of teachers for these events
\end{itemize}

\section{Evaluation}
\subsection{Execution Time}
\begin{itemize}
	\item comparison of the different execution times - only checking satisfiability with Z3 and optimizing with PuLP, vZ and the binary implementation (with various cut-off points)
	\item conclusion: PuLP much faster than vZ and also for Z3 (if using realistic cutoff points), checking for satisfiability very fast
	\item $\rightarrow$ PuLP is used in the final version of the program
\end{itemize}

\subsection{Output Quality}
\begin{itemize}
	\item the calculated solutions are sensible and in some ways better than the solutions obtained by hand
	\item it becomes apparent that lots of bias is included when assignment process is done by hand - intended or unintended
	\item many more soft objectives have to be included, e.g. teaching quality of teachers, to take all of the intended bias into account
	\item making it possible for the end-user to pre-enter teachers for certain events  and to change weights, makes it possible to find a good solution in little time, even with very specific objectives/constraints
\end{itemize}



\section{Conclusion}
\begin{itemize}
	\item is the problem solved?
	\item what are the advantages / disadvantages of this approach
	\item which solvers proved to be useful
	
\end{itemize}






\end{document}
